\section{Marco de problemas inversos}

La idea central de la segunda categoría de técnicas de guía es modelar la generación condicional como un \textbf{problema inverso}. Se trata de un marco más estructurado y general, aplicable a una gama más amplia de escenarios.

\subsection{Definición del problema inverso}

\begin{importantbox}{Definición formal del problema inverso}
Dado:
\begin{itemize}
    \item Una función de mapeo hacia adelante \textbf{forward mapping function} $\mathcal{A}$, que mapea datos limpios $x$ a observaciones parciales $y$: $y = \mathcal{A}(x) + n$, donde $n$ es el posible ruido de observación.
    \item Un valor de observación $y$.
\end{itemize}
Objetivo: encontrar un $x$ razonable tal que $\mathcal{A}(x) \approx y$.

Dado que $\mathcal{A}$ suele ser \textbf{no invertible}, pueden existir múltiples $x$ razonables para un $y$ dado; por lo tanto, esto es esencialmente un problema de generación condicional.
\end{importantbox}

\subsection{Aplicaciones típicas del problema inverso}

Los problemas inversos tienen numerosas aplicaciones en el procesamiento de imágenes y pueden extenderse a otras modalidades de datos:

\begin{table}[H]
\centering
\begin{tabular}{lll}
\toprule
\textbf{Tarea} & \textbf{Mapeo hacia adelante $\mathcal{A}$} & \textbf{Descripción} \\
\midrule
Denoising gaussiano (desenfoque) & Núcleo de convolución gaussiano & Recuperar una imagen nítida desde una borrosa \\
Superresolución & Operación de submuestreo & Recuperar alta resolución desde baja resolución \\
Inpainting de imágenes & Operación de máscara & Recuperar una imagen completa desde observaciones parciales \\
Reconstrucción CT/MRI & Operación de proyección & Reconstruir 3D desde proyecciones 2D multángulo \\
Colorización de imágenes & Operación de descoloración & Recuperar una imagen a color desde un mapa de grises \\
\bottomrule
\end{tabular}
\caption{Escenarios típicos de aplicación del problema inverso}
\end{table}

\begin{knowledgebox}{Generalidad del problema inverso}
Los problemas inversos no se limitan al dominio de las imágenes. En áreas como el diseño molecular, el plegamiento de proteínas y la planificación robótica, cualquier problema que implique ``inferir información completa desde observaciones parciales'' puede modelarse como un problema inverso. Si puedes definir un mapeo hacia adelante $\mathcal{A}$ razonable para un nuevo dominio, podrás utilizar las técnicas presentadas en esta clase para resolverlo.
\end{knowledgebox}

\subsection{Resumen de la sección}

El problema inverso proporciona un marco unificado que conecta la guía de generación durante la inferencia con los clásicos problemas de procesamiento de señales/imagen computacional. Los elementos clave son: el mapeo hacia adelante conocido $\mathcal{A}$, el valor de observación $y$ y el modelo de generación incondicional preentrenado. A continuación, derivaremos cómo utilizar la fórmula de Bayes para transformar la solución del problema inverso en una modificación de funciones puntuales (score functions).

